\section{Specialness on every upper model}\label{sec:specialness}

We prove part~(i) of Theorem~\ref{thm:example} by excluding every
positive-dimensional general-type orbifold base on every allowed upper
model. For surface bases we use the growth of logarithmic
pluricanonical sections. For curve bases we use the complete elliptic
fibres of $j$ and its retained multiplicity-one components.

Let
\[
 u:(Y',D')\longrightarrow(Y,D)
\]
be any elementary birational orbifold morphism, with $(Y',D')$ smooth
projective, and let $g:(Y',D')\to Z$ be any positive-dimensional neat
fibration. Every component of $D'$ is either the strict transform of one
of the $D_i$ or a $u$-exceptional curve, because $u_*D'=D$. The images of
the exceptional curves form a finite subset of $Y$. There are two
possibilities for the dimension of $Z$.

\subsection{A two-dimensional base}
Suppose $\dim Z=2$. A general connected zero-dimensional fibre is a point,
so $g$ is birational. Compare $g$ with the identity fibration on
$(Y',D')$, using the identity on the source and $g$ on the base. By the
definition of canonical dimension,
\begin{equation}\label{eq:birational-kappa}
 \kappa(g,D')\leq\kappa(Y',K_{Y'}+D').
\end{equation}

We claim that the right side is at most one. Choose a positive integer
$m$ divisible enough that $mD'$ is integral. Off the finite set of
exceptional images, $u$ is an isomorphism and its pushed-forward
boundary is $D$. A section of $m(K_{Y'}+D')$ therefore restricts to a
section of $m(K_Y+D)$ away from that finite set. A section of a line
bundle on a smooth surface extends across a subset of codimension two:
after trivializing the line bundle, this is extension of regular
functions on a normal variety. We obtain an injection
\begin{equation}\label{eq:section-injection}
 H^0\bigl(Y',m(K_{Y'}+D')\bigr)
   \lhook\joinrel\longrightarrow
 H^0\bigl(Y,m(K_Y+D)\bigr).
\end{equation}
The target grows linearly with $m$ by \eqref{eq:log-canonical}. This proves
the claim. Neatness, \eqref{eq:neat-kappa}, and
\eqref{eq:birational-kappa} now show
\[
 \kappa\bigl(Z,K_Z+\Delta(g,D')\bigr)\leq1<2.
\]
Notice that this uses a one-sided bound on sections; it does not assume
birational invariance for arbitrary changes of the boundary.

\subsection{A curve base}
Suppose $Z=C$ is a smooth projective curve, and set
$\Delta=\Delta(g,D')$. Assume for contradiction that
\begin{equation}\label{eq:positive-base}
 \deg(K_C+\Delta)>0.
\end{equation}
Choose $t\in\PP^1$ away from the $q_i$ and from the images, under $j$, of
the finitely many exceptional images of $u$. The fibre
\[
 C_t=(j\circ u)^{-1}(t)
\]
is then an unchanged smooth complete elliptic curve, disjoint from both
$D'$ and the exceptional locus of $u$.

Suppose $h=g|_{C_t}:C_t\to C$ is nonconstant. It is finite and
surjective. For a point $c\in C$, write $m_c$ for its orbifold
multiplicity. Any component $P$ of $g^*c$ meeting $C_t$ lies outside
$\Supp D'$, so $m_P(D')=1$. Formula~\eqref{eq:infimum} implies
\[
 r(g,P)\geq m_c.
\]
For $x\in C_t$ above $c$, the ramification index $e_x$ of $h$ is the
intersection multiplicity of $C_t$ with $g^*c$ at $x$. It is the sum of
the positive intersection orders multiplied by the corresponding
$r(g,P)$, so $e_x\geq m_c$. If $m_c=\infty$, this is impossible, since
surjectivity supplies a point $x$ above $c$. Thus the existence of $h$
first rules out all infinite base multiplicities. For finite $m_c$,
\[
 e_x-1\geq e_x\left(1-\frac1{m_c}\right).
\]
Summing this inequality over the fibres above $\Supp\Delta$ and using
Riemann--Hurwitz gives
\begin{align*}
 0=\deg K_{C_t}
 &=\deg(h)\deg K_C+\sum_{x\in C_t}(e_x-1)\\
 &\geq\deg(h)\deg(K_C+\Delta)>0,
\end{align*}
a contradiction. This argument only needs the inequality of
multiplicities, so it applies to the infimum convention rather than
requiring divisibility of ramification indices.

It follows that $g$ is constant on all sufficiently general fibres of
$j\circ u$. Those fibres are connected. Descent on function fields gives
a rational map $r:\PP^1\dashrightarrow C$ with
\[
 g=r\circ j\circ u.
\]
A rational map between smooth projective curves extends everywhere. If
$\deg r>1$, a general fibre of $g$ is the disjoint union of the fibres of
$j\circ u$ over the distinct preimages of a general point of $C$; this
contradicts connectedness of the general fibre of $g$. Therefore
$\deg r=1$, and $C\simeq\PP^1$. Up to this isomorphism, $g=j\circ u$.

It remains to compute the boundary of this possible curve base on the
arbitrary model $Y'$. Let $A_i'$ be the strict transform of $A_i$.
Its coefficient in $D'$ is zero, since $u_*D'=D$ and $A_i$ does not
occur in $D$. At its generic point $u$ is an isomorphism. Thus $A_i'$
occurs with multiplicity one in $(j\circ u)^*q_i$, and
\[
 m_{j\circ u}(q_i)=1.
\]
For every $t\notin\{q_i\}$ the strict transform of the original elliptic
fibre similarly has boundary coefficient zero and ordinary multiplicity
one. Additional exceptional components cannot increase an infimum
already equal to one. Consequently
\[
 \Delta(j\circ u,D')=0.
\]
The degree in \eqref{eq:positive-base} is therefore $\deg K_{\PP^1}=-2$,
the final contradiction. Both possible base dimensions have been
excluded, proving specialness on every upper model in
Definition~\ref{def:special}.

\begin{remark}
This proof is compatible with adding boundary components on further
exceptional curves allowed by the definition. It never adds the
strict transforms of the original $A_i$ to the boundary: their
coefficients must remain zero because their images are divisors absent
from $D$. This is why the universal quantifier over upper models does
not remove the example.
\end{remark}
